% Independent contribution by /root/frobenius.
% Uses citation keys from cycles/references.bib and C4_references.bib.
\subsection{C.4 supplement: exact integral bounds without a principal polarization}
\label{sec:c4-nonprincipal-supplement}

\paragraph{C.4h (Proved: the ideal of scalar factorizations).}
Let $A$ be an abelian variety of positive dimension $g$ over an
algebraically closed field $k$, in arbitrary characteristic. We define
\[
 I_J(A)=\{n\in\mathbb Z:\ [n]_A=v\circ u
       \text{ for some }A\xrightarrow{u}\prod_jJ(C_j)
                           \xrightarrow{v}A\},
\]
where each $C_j$ is a smooth projective connected curve over $k$.
This is an ideal of $\mathbb Z$. Indeed, multiplication of $v$ by an
integer realizes integer multiples; the product of two intermediate
products of Jacobians realizes the sum of their two composites.
We will prove below that $I_J(A)\ne0$, and write
$I_J(A)=s_J(A)\mathbb Z$ with $s_J(A)>0$.
Then
\begin{equation}\label{eq:c4-exact-splitting}
 A\text{ is an actual direct summand of a product of Jacobians}
 \quad\Longleftrightarrow\quad s_J(A)=1.
\end{equation}
For completeness, if $v u=1_A$, the maps
\[
 A\times\ker v\longrightarrow J,\quad(a,b)\longmapsto u(a)+b,
 \qquad
 J\longrightarrow A\times\ker v,\quad
 x\longmapsto(v(x),x-u v(x))
\]
are mutually inverse as group schemes. The kernel is a connected smooth
proper group: smoothness follows from the splitting of the differential
of $v$, and connectedness from the displayed product decomposition.
It is consequently an abelian variety.

The definition also gives
\[
 s_J(\widehat A)=s_J(A),\qquad
 s_J(A\times B)=\operatorname{lcm}(s_J(A),s_J(B)).
\]
For the first assertion we dualize a factorization and use the canonical
principal polarizations of the Jacobians. For the second we restrict and
project a factorization to each factor, and conversely combine two
factorizations of the same integer.

Here is the exact isogeny calculation needed later. If
$f:A\to B$ is an isogeny and the \emph{group scheme} $\ker f$ is killed
by $e>0$, the quotient property gives a homomorphism $h:B\to A$ with
$h f=[e]_A$; surjectivity of $f$ also gives $f h=[e]_B$.
Thus
\begin{equation}\label{eq:c4-isogeny-divisibility}
 eI_J(B)\subset I_J(A),\qquad eI_J(A)\subset I_J(B).
\end{equation}
In particular an isogeny whose kernel is killed by an integer prime to
$\ell$ preserves the $\ell$-adic valuation of $s_J$.
We may always take $e=\deg f$; this remains true for inseparable
isogenies \citep[Section 8, paragraph after Proposition 8.1,
p.~12 of the corrected text]{frC4Milne1986}.
None of these statements replaces an integral retraction by a
rational inverse to an isogeny.

\paragraph{C.4i (Proved: a uniform factorial bound in every characteristic).}
For every $A$ as above,
\begin{equation}\label{eq:c4-factorial-bound}
             s_J(A)\mid(2g-1)!.
\end{equation}
We spell out the algebraic construction, including the integer in the
composite. This is a quantitative application of the proof of
\citet[Proposition 3.4, equations (3.1)--(3.3)]{cyScaviaSuzuki2025},
not an inference from rational cohomology alone.

We first record its useful form. Suppose $X$ has dimension $d$,
$\theta=c_1(L)\in\mathrm{NS}(X)$, and an integral curve cycle
$\alpha$ has numerical class
\[
                  \alpha=n\,\frac{\theta^{d-1}}{(d-1)!}.
\]
Write $\alpha=\sum_j m_j(f_j)_*[C_j]$ using normalizations of its
integral curve components. Such normalizations are smooth over the
algebraically closed field, including in characteristics $2$ and $3$.
Choose base points and let $\bar f_j:J(C_j)\to X$ be the induced
homomorphisms. Set $J=\prod_jJ(C_j)$, let
$\lambda_j:J(C_j)\xrightarrow{\sim}\widehat{J(C_j)}$ be its canonical
polarization, and define
\[
 U:X\longrightarrow J,\quad
 U=(\lambda_j^{-1}\widehat{\bar f_j}\,\varphi_\theta)_j,
 \qquad
 V:J\longrightarrow X,\quad V=\sum_jm_j\bar f_j.
\]
Matsusaka's identity gives the equality of homomorphisms
\begin{equation}\label{eq:c4-matsusaka-integer}
               VU=[\,n\chi(L)\,]_X,\qquad
               \chi(L)=\deg\frac{\theta^d}{d!}.
\end{equation}
The sign is the positive sign in equation (3.1) of the cited proof:
its two occurrences of Matsusaka's $\delta$ have the same negative
sign and cancel. Negative $m_j$ are carried by $V$, and cause no
positivity requirement on $\theta$. Translating $f_j$ to send its
base point to zero does not change its numerical curve class.

We can also check the coefficient in \eqref{eq:c4-matsusaka-integer}
directly. In an $\ell$-adic symplectic basis for the nondegenerate
alternating form of $\theta$, write the form as
$\sum_{a=1}^{d}t_a e_a\wedge f_a$.
The complementary exterior power
$\theta^{d-1}/(d-1)!$ has coefficients
$\prod_{b\ne a}t_b$. Multiplication by the original form therefore
gives the scalar $\prod_b t_b=\chi(L)$.
The same coefficient follows for every $\theta$ by a polynomial
identity. Faithfulness of the Tate module realization of
homomorphisms of abelian varieties makes this an equality of
homomorphisms, as in the cited Matsusaka proof.

Now take $X=A\times\widehat A$, $d=2g$, and
$\theta=c_1(\mathcal P_A)$. The associated symmetric homomorphism
$\varphi_\theta:X\to\widehat X=\widehat A\times A$ exchanges the two
factors, so it is an isomorphism. The Riemann--Roch formula
$\deg\varphi_\theta=\chi(\mathcal P_A)^2$ gives
$\chi(\mathcal P_A)=\pm1$ (in the usual orientation it is $(-1)^g$).
The integral intersection cycle $\alpha=\theta^{2g-1}$ satisfies the
preceding hypothesis with $n=(2g-1)!$.
Equation \eqref{eq:c4-matsusaka-integer} thus factors multiplication
by $\pm(2g-1)!$ on $X$ through $J$.
We adjust the sign of $V$, then compose with the inclusion of $A$
and the projection to $A$. This proves \eqref{eq:c4-factorial-bound}.
The Riemann--Roch and Poincar\'e degree facts are also explicitly
used in the proof of Corollary 1.3 of
\citet{cyScaviaSuzuki2025}.

\paragraph{C.4j (Proved: a simultaneous polarization criterion in dimensions at most three).}
Define the positive integer
\[
 \delta(A)=\gcd\{\chi(L):L\in\mathrm{Pic}(A)\}.
\]
Equivalently, we may take the gcd only over ample $L$.
To see this, fix an ample $H$ and any line bundle $L$.
The function $t\mapsto\chi(L\otimes H^t)$ is a polynomial of degree at
most $g$ with integer values. For all sufficiently large consecutive
integers $t$, the line bundles are ample. Integral finite differences,
or interpolation in the binomial basis shifted to these integers,
express its value at $0$ as an integral linear combination of these
values. Hence their common divisor divides $\chi(L)$.

If $1\le g\le3$, then
\begin{equation}\label{eq:c4-polarization-gcd}
 \boxed{\quad s_J(A)\mid
             \gcd\bigl((2g-1)!,\delta(A)\bigr).\quad}
\end{equation}
Indeed, let $L$ be ample. The group-scheme form of the polarization
descent theorem gives a maximal isotropic subgroup scheme
$H\subset K(L)$, a quotient isogeny $f:A\to B=A/H$, and a principal
polarization $M$ on $B$ with $f^*M\simeq L$.
This is \citet[Section 23, Theorem 4 and Corollary 1,
pp.~233--234]{frC4Mumford1970}; its proof explicitly includes
$\mu_p$ and $\alpha_p$, so it has no separability restriction.
Riemann--Roch gives
\[
                       \deg f=\chi(L).
\]
By \citet[Theorem 4, p.~2 of the opened preprint]{cyOortUeno1973}, the principally
polarized variety $B$, in these dimensions and over any algebraically
closed field, is the Jacobian of a compact-type curve, and therefore
the product of the Jacobians of its smooth irreducible components.
Thus $\chi(L)\in I_J(A)$ by \eqref{eq:c4-isogeny-divisibility}.
Taking the gcd and using \eqref{eq:c4-factorial-bound} proves the claim.

Consequently it suffices that, for \emph{each} prime
$\ell\mid(2g-1)!$, there exist an ample line bundle $L_\ell$ on $A$
with $\ell\nmid\chi(L_\ell)$.
The choices of $L_\ell$ may be different. We take the finitely many
resulting factorizations, together with the factorial factorization,
and integers $a_j$ with $\sum_j a_j n_j=1$. If
$v_j u_j=[n_j]_A$, then
\[
                 u=(u_j)_j,\qquad v=\sum_j a_jv_j
                 \quad\text{satisfy}\quad vu=1_A.
\]
This constructs a single actual retraction. In particular:
\[
\begin{array}{c|c|c}
 \dim A & \text{universal bound for }s_J(A)
        &\text{sufficient local polarization checks}\\ \hline
 2 & 6 & \ell=2,3\\
 3 & 120 & \ell=2,3,5.
\end{array}
\]
If $k$ has characteristic $2$ or $3$, the check at that very prime
is retained. The construction uses finite subgroup schemes rather
than their sets of geometric points.

The kernel exponent can improve the quantitative divisibility:
if a principalization has kernel killed by $e$, we obtain
$s_J(A)\mid\gcd((2g-1)!,e)$.
For a separable polarization of type
$(d_1,\ldots,d_g)$, its kernel has exponent $d_g$, so we may use
$e=d_g$. For an inseparable polarization we use its actual
group-scheme kernel exponent, or the always valid bound
$e=\chi(L)$; a prime-to-characteristic elementary-divisor
description must not be substituted for the full kernel.

\paragraph{C.4k (Proved: what the N\'eron--Severi calculation establishes).}
For an abelian surface, choose a basis $D_1,\ldots,D_\rho$ of
$\mathrm{NS}(A)$, and write
$q_i=D_i^2/2$ and $b_{ij}=D_iD_j$. Then
\begin{equation}\label{eq:c4-delta-surface}
 \delta(A)=\gcd\bigl(\{q_i\}_i,\{b_{ij}\}_{i<j}\bigr).
\end{equation}
Indeed,
$\chi(\sum n_iD_i)=\sum_iq_in_i^2+\sum_{i<j}b_{ij}n_in_j$;
the reverse divisibility follows by evaluating at $D_i$ and
$D_i+D_j$. Thus this particular sufficient test is a finite
integral lattice calculation.
In dimension three we may compute the analogous gcd from the
Euler polynomial on the finite grid $0\le n_i\le3$: its successive
finite differences are integral, and the binomial-basis
interpolation expresses every other integer value as an integral
combination of these grid values.

If $\mathrm{NS}(A)=\mathbb Z[L]$ for a primitive ample $L$, then
$\delta(A)=\chi(L)$, since $\chi(L^n)=n^g\chi(L)$.
In particular a complex surface with primitive polarization type
$(1,d)$ is settled whenever $\gcd(d,6)=1$; a complex threefold
with primitive polarization type $(1,d_2,d_3)$ is settled whenever
$\gcd(d_2d_3,30)=1$.
These are sufficient conditions and give no counterexample when
the gcd is larger than one.

There is a precise limitation on the proposed coprime-isogeny
method. An isogeny $A\to B$ of degree prime to $\ell$, with $B$
principally polarized, pulls a principal line bundle back to an
ample $L$ with $\ell\nmid\chi(L)$. The converse follows from
Mumford's principalization. Therefore, when $\ell\mid\delta(A)$,
there is \emph{no} such prime-to-$\ell$ principalization in the same
dimension, even if we allow all possible quotient polarizations.
This rules out that particular proof strategy in the residual
Picard-rank-one cases; it does not rule out a retraction through
Jacobians of higher dimension.

\paragraph{C.4l (Proved limitation of divisor products; no geometric counterexample asserted).}
We can make the residual surface calculation more explicit.
Let $A/\mathbb C$ have $\mathrm{End}(A)=\mathbb Z$ and
$\mathrm{NS}(A)=\mathbb Z[L]$, with $L$ of type $(1,d)$.
Let $\widehat L$ be the primitive dual polarization, normalized by
$\varphi_{\widehat L}\varphi_L=[d]$.
On $A\times\widehat A$ put
$D=\operatorname{pr}_1^*c_1(L)$,
$\widehat D=\operatorname{pr}_2^*c_1(\widehat L)$ and
$P=c_1(\mathcal P_A)$.
The N\'eron--Severi group is generated by
$D,\widehat D,P$: the cross summand is
$\operatorname{Hom}(A,\widehat{\widehat A})=\operatorname{End}(A)$.
The endomorphisms induced on $\widehat A$ by integral curve cycles
that are integer combinations of products of three divisors form
exactly
\begin{equation}\label{eq:c4-divisor-product-ideal}
                         \gcd(6,d)\,\mathbb Z\,1_{\widehat A}.
\end{equation}

Here is a direct calculation, avoiding an assumption that all
integral curve classes are generated by divisors. Choose an
integral symplectic basis $x_1,x_2,x_3,x_4$ for $H^1(A,\mathbb Z)$
and the dual basis $y_1,\ldots,y_4$ for $H^1(\widehat A,\mathbb Z)$
so that
\[
 D=x_1x_2+d\,x_3x_4,\quad
 \widehat D=d\,y_1y_2+y_3y_4,\quad
 P=\sum_{a=1}^{4}x_ay_a.
\]
Products here mean exterior products. The normalized class
$P^3/6$ has the cross component inducing $1_{\widehat A}$, and
expansion gives
\begin{equation}\label{eq:c4-surface-exterior}
                 D\widehat D P=-d\,P^3/6.
\end{equation}
The only degree-three monomials in these divisor classes with
K\"unneth bidegree $(3,3)$ are $P^3$ and $D\widehat D P$;
the other monomials have no component inducing a map
$\widehat A\to\widehat A$. They therefore give precisely the
integers generated by $6$ and $d$.
The identification of the cross component follows either from
the curve-factorization formula in C.4a or from
\eqref{eq:c4-matsusaka-integer} for $P$.
The script \texttt{verify\_c4\_exterior.py} checks the exterior
identity symbolically, together with its factorial normalization.

Thus for $d=2$, $3$, or $6$ the divisor-product calculation itself
cannot produce the desired identity. This is a proved statement
about that subring of algebraic cycles, not a statement about
all of $\mathrm{CH}_1(A\times\widehat A)$.
The rational Hodge group calculation does not justify replacing
the latter integral group by the former.

\paragraph{The scope of the available comparison theorems.}
The opened results impose the following precise limits.
The converse from the integral Hodge conjecture for curves on
$A$ to the direct-summand property requires a unimodular line
bundle in \citet[Theorem 1.1 and the paragraph immediately
following it]{cyBeckmannFortman2023} and
\citet[Theorem 1.1 and the paragraph following it]{cyScaviaSuzuki2025}.
Without that hypothesis the applicable complex criterion is on
$A\times\widehat A$
\citep[Theorem 1.4]{cyVoisin2024}.
Zarhin's construction makes $(A\times\widehat A)^4$ principally
polarizable, but its dimension is $8g$, beyond the
dimension-three Jacobian theorem; it alone supplies no retraction
through Jacobians \citep[Remarks 16.12--16.13]{frC4Milne1986}.
Finally the opened 2026 counterexample theorem concerns very
general principally polarized varieties of dimension at least
four \citep[Theorem 1.1 and Corollary 1.5]{cyEngelFortmanSchreieder2026}.
Its hypotheses do not include a product $A\times\widehat A$
with $A$ a nonprincipal surface. Applying its degeneration
obstruction to that restricted family would require a new argument.
